\section{Related Work}
\label{sec:related}

\paragraph{Food physics simulation and manipulation.}
Physics-based simulation of food materials spans robotics and computer graphics. Ding~et~al.~\cite{ding2019thermomechanical} introduced a thermomechanical MPM for simulating baking. DiSECt~\cite{heiden2021disect} proposed a differentiable simulator for cutting and used Bayesian inference to calibrate material parameters from force measurements. RoboNinja~\cite{xu2023roboninja} trained adaptive cutting policies with differentiable MPM, and RoboCook~\cite{shi2023robocook} addressed long-horizon elasto-plastic manipulation. CulinaryCut-VLAP~\cite{koh2026culinarycut} coupled vision-language models with MPM for food cutting. Most closely related to our inverse estimation objective, SliceIt!~\cite{beltran2024sliceit} calibrated food material parameters from real-world force measurements in a dual-simulator pipeline for robotic slicing. On the visual side, FruitNinja~\cite{wu2025fruitninja} generated realistic interior textures for 3D Gaussian Splatting fruit models, and PhysGaussian~\cite{xie2024physgaussian} integrated MPM with Gaussian Splatting for physics-driven dynamics. These works either assume differentiable simulators, calibrate from direct force measurements, or specify parameters manually. We address the inverse problem of discovering material parameters for a \emph{non-differentiable} fracture simulator from behavioral feature targets alone.

\paragraph{Material Point Method, fracture, and inverse problems.}
The Material Point Method combines Lagrangian particles with an Eulerian grid to handle large deformations and fracture. Wolper~et~al.~\cite{wolper2019cdmpm} introduced continuum damage mechanics into MPM, and AnisoMPM~\cite{wolper2020anisompm} extended this to anisotropic damage with directional fiber reinforcement. These simulators are non-differentiable: topological discontinuities during fracture preclude gradient computation through the simulation. This distinguishes our setting from differentiable physics frameworks such as DiffTaichi~\cite{hu2020difftaichi}, FluidLab~\cite{xian2023fluidlab}, and PlasticineLab~\cite{huang2021plasticinelab}, where gradients enable direct inverse solving. NCLaw~\cite{ma2023nclaw} demonstrated that neural constitutive laws can be learned within a differentiable MPM, but this requires backpropagation through the simulator, which is fundamentally unavailable for fracture mechanics.

\paragraph{Latent space optimization and reinforcement learning.}
Optimization in learned latent spaces was pioneered by G\'{o}mez-Bombarelli~et~al.~\cite{gomez2018automatic} for molecular design, followed by weighted retraining~\cite{tripp2020sample} to focus generative models on high-scoring regions. LOL-BO~\cite{maus2022lolbo} added trust regions to latent Bayesian optimization, and Lee~et~al.~\cite{lee2025nfbo} identified a value discrepancy in VAE-based optimization that normalizing flows eliminate through bijective mappings. In robotics, PULSE~\cite{luo2024pulse} demonstrated that distilling motor skills into a latent space enables efficient hierarchical RL, and OmniGrasp~\cite{luo2024omnigrasp} applied this for dexterous grasping. We adapt the latent-space-for-control principle from the motor domain to parameter identification.

\paragraph{Surrogate-based optimization and amortized inference.}
Training a neural surrogate on a pre-computed dataset and optimizing against it is formalized as offline model-based optimization (MBO), benchmarked through Design-Bench~\cite{trabucco2022designbench}. Surrogate exploitation, where the optimizer discovers inputs scoring highly on the surrogate but poorly under the true objective, is a well-known failure mode addressed by conservative training~\cite{trabucco2021coms} and robustness constraints~\cite{yu2021roma}. CMA-ES~\cite{hansen2001cmaes} is the standard gradient-free optimizer in this setting~\cite{fischer2024zerograds}. Amortized inference trains a single model to solve a distribution of inverse problems~\cite{cranmer2020frontier, amos2023tutorial}. Ardizzone~et~al.~\cite{ardizzone2019inn} used invertible neural networks for physics inverse problems. Goal-conditioned policies~\cite{schaul2015universal} generalize across objectives by conditioning on a target specification.
